\chapter{绪论}
\renewcommand{\currentchapterpath}{chapters/ch01/}

\begin{flushright}
\textit{“知止而后有定，定而后能静，静而后能安，安而后能虑，虑而后能得。” ——《大学》
\footnote{强调先明确目标与边界，再逐步深入思考。学习多体理论同样需要先把握问题的尺度与语言。}}
\end{flushright}

量子多体理论（quantum many-body theory）研究由大量微观粒子组成的量子体系：电子、声子、自旋、冷原子等。与少数体问题不同，粒子数 $N$ 通常达到 $10^{23}$ 量级，精确对角化既不可能也无必要；关键在于识别集体自由度、建立有效描述，并从可观测量出发理解宏观物性。

本章给出问题的整体图景与历史线索，为后续二次量子化、格林函数、平均场/密度泛函以及线性响应等章节做准备。


\chapterinput{overview}
\chapterinput{history}
